\SourceLeaf{9}
\Article{2}
\par
We will consider, in the greatest generality, a finite body of any shape, separated from the rest of infinite space by one continuous surface or by several continuous surfaces distinct from one another (if the body happens to enclose one or more cavities). We will call the totality of these surfaces simply the surface of the body. Imagine this surface divided into an infinity of elements $\mathrm{d}s$. Let $P$ be a point of the element $\mathrm{d}s$, and let its coordinates referred to three mutually perpendicular planes be denoted by $x,y,z$. Let $PX,PY,PZ$ be straight lines parallel to the respective coordinate axes and directed toward the sides on which the coordinates are assumed to take positive increments. Let $PQ$ further be the outward-directed normal to the surface. Let $M$ be the attracted point, situated anywhere, let its coordinates be $a,b,c$, and let the distance $PM$ (always to be taken as positive) be $=r$. We will denote the angles that the line $PM$ makes with $PX,PY,PZ$ by $MX,MY,MZ$, and the angles between $PQ$ and $PX,PY,PZ,PM$ by $QX,QY,QZ,QM$. All this refers to points of the surface without specifying which points. Whenever we need to deal with several definite points of the surface, we will use the same symbols distinguished by primes.

\Article{3}
\par
Imagine a plane perpendicular to the $x$-coordinate axis, situated so that, if its equation is written as $x=\alpha$, $\alpha$ is less than the least value of the coordinate $x$ on the surface of the body. The body projected onto this plane will trace a finite figure there, which we suppose divided into infinitely many elements $\mathrm{d}\Sigma$. In the element $\mathrm{d}\Sigma$, erect a perpendicular at a point $\Pi$ (thus parallel to the $x$-coordinate axis) that intersects the body at the points $P',P'',P'''$ etc. The number of these points will evidently be even. Also erect perpendiculars to the plane at every point on the boundary of the element $\mathrm{d}\Sigma$. They will form a cylindrical surface in the broader sense and cut out elements $\mathrm{d}s',\mathrm{d}s'',\mathrm{d}s'''$ etc. from the surface of the body. The element $\mathrm{d}\Sigma$ will be the projection of each of the elements $\mathrm{d}s',\mathrm{d}s'',\mathrm{d}s'''$ etc. Hence $\mathrm{d}\Sigma=\pm\mathrm{d}s'.\cos QX'=\pm\mathrm{d}s''.\cos QX''=\pm\mathrm{d}s'''.\cos QX'''$ etc., with the upper or lower sign applying according as the cosine is that of an acute or an obtuse angle. Since the perpendicular evidently enters the body at $P'$, leaves it at $P''$, enters it again at $P'''$, and so on, it is easy to see that $QX'$ is obtuse, $QX''$ acute, $QX'''$ obtuse, and so on, so that we have
